\scanpage{40}\mspage{39}
% L1–3
\uline{NB}\qquad Soit $(X_\alpha)$ famille \inserted{finie} de parties modérées de $L$, \deleted{mut\textsuperscript{t}.} mutuell\textsuperscript{t} modérées i.e telles que
% L4
\[
S=\bigcup_\alpha X_\alpha^{\bullet}
\]
% L5–6
soit une partie tot ordonnée~:
\[
t_1<t_2<\ldots<t_n
\]
% L7–12
Alors l'Alg. de Boole engendrée par les $X_\alpha$ \deleted{est \unsure{consti|enti}} admet comme \deleted{composantes} \deleted{\unsure{éléments} du spectre (quand on les interprète comme parties de $L$,} ``composants'' les ens modérés \inserted{(et ps.tronçons)} suivants (rangés par ordre croissant pour $<$)
% L13–14
\[
\scriptstyle
\underset{\substack{\vert\\\inserted{\text{si }t_1\notin\partial L}}}{[L_{<t_1}]}\unsure{,|a}\;
\{t_1\},\ ]t_1,t_2[,\ \{t_2\},\ ]t_{1},t_3[,\ \{t_3\},\ldots,\
\deleted{\unsure{\{t\}}}\ \{t_n\},\
\underset{\substack{\rule{3em}{0.35pt}\\\inserted{\text{si }t_n\notin\partial L}}}{[L_{>t_n}]}
\]
% L15–18
Ainsi, cette alg. de Boole est connue quand on connait $S$, i.e. les points $t_1<\ldots<t_n$ (\unsure{ou la|ou} ``décomposition'' \inserted{ou ``subdivision''} du tronçon \unsure{$L$|$\Lambda$|$A$}).
